\documentclass[conference]{IEEEtran}
\usepackage{amsmath,amssymb,amsfonts}
\usepackage{booktabs}
\usepackage{graphicx}
\usepackage{textcomp}
\usepackage{xcolor}
\usepackage[hidelinks]{hyperref}
\usepackage{multirow}
\usepackage{array}

\begin{document}

\title{Component-Wise Quantization of Vision--Language--Action Policies:\\
Activations Are the Wall, Rotations Are the Key}

\author{
\IEEEauthorblockN{Justin Williams, Kishor Datta Gupta, Roy George, Mrinmoy Sarkar}
\IEEEauthorblockA{Clark Atlanta University, Atlanta, GA\\
\texttt{justinwilliamstech693@gmail.com}}
}

\maketitle

\begin{abstract}
Vision--Language--Action (VLA) policies built on 7B-parameter language backbones are accurate
but too large for on-board robot inference. We present a comprehensive, closed-loop
post-training quantization (PTQ) study of an OpenVLA-OFT policy on the LIBERO benchmark,
covering the full quantization design space from mechanism diagnosis to component-wise
allocation to real edge-hardware confirmation. Our main findings: \textbf{(1) The bottleneck
is activations, not weights}---naive 4-bit weight quantization is near-lossless (90.75\% vs.\
a 90.25\% bf16 baseline), but naive 4-bit activation quantization collapses the policy to 0\%
because a single backbone layer carries a channel $10^5\times$ larger than the median, forcing a
quantization scale that zeros all other activations. \textbf{(2) A single orthogonal rotation
fixes it}---applying $Wx=(WQ^\top)(Qx)$ before quantization spreads that outlier energy
uniformly, cutting the median activation outlier ratio from 18.0 to 1.6 and restoring W4A4
to 88.8--90.5\% ($\geq$98\% retention). We benchmark \textbf{16 transforms}; all global
dense mixers work ($\sim$27\% A4 error); PCA is actively harmful (it concentrates energy);
the DCT is the deployment pick (data-free, $O(n\log n)$, no power-of-two fallback).
\textbf{(3) Component-wise allocation reveals that the vision tower is free (INT8/FP8 both
lossless) and the L1-regression action head is the floor (INT8 safe, INT4 collapses to 70\%)}.
Composing these: INT8 vision $+$ DCT-W3A8 backbone $+$ INT8 head achieves
\textbf{74\% compression at 89.5\% success vs.\ 88.0\% bf16---no measurable accuracy loss} at
400 rollouts. \textbf{(4) We confirm the result on real Jetson AGX Orin hardware}: W4A4+DCT
achieves 86.0\% on LIBERO-Spatial ($-$2.2 pt vs.\ A100 bf16; 97.5\% retention;
statistically indistinguishable at 100 rollouts)---the first closed-loop real-hardware
confirmation of a quantized L1-regression VLA. Code and harness are released.
\end{abstract}

\begin{IEEEkeywords}
VLA quantization, activation outliers, incoherence rotation, DCT, component-wise PTQ,
edge deployment, Jetson Orin, LIBERO
\end{IEEEkeywords}

% ============================================================================
\section{Introduction}
\label{sec:intro}
% ============================================================================

Vision--Language--Action models adapt large vision--language backbones into robot policies
that map an image and a language instruction to continuous actions~\cite{openvla,openvlaoft}.
Their accuracy rests on 7B-parameter backbones (Llama-2-7B), which demand 15+ GB of memory
and >300\,W GPU power---precluding deployment on the vehicle-class compute ($\leq$60\,W) a
mobile robot can carry. Post-training quantization (PTQ) compresses such models without
retraining, but most PTQ research optimizes language-model perplexity or open-loop token
accuracy~\cite{gptq,awq,smoothquant,quarot,spinquant}, not the \emph{closed-loop task success}
where per-step errors compound over rollouts of 200--600 simulator steps.

This paper is an end-to-end study, structured around three coupled questions:
\emph{(i) Where does quantization hurt a VLA policy?}
\emph{(ii) What is the cheapest fix, and which variant of the fix is best?}
\emph{(iii) How far can one compress, and what does the compressed model look like on real
edge hardware?} Our contributions:

\begin{enumerate}
  \item \textbf{Mechanism diagnosis.} Through a calibrated outlier ablation, we identify
        the specific layer (\texttt{layers.1.mlp.down\_proj}) and the specific failure mode
        (one channel $\sim\!10^5\times$ the median) that causes W4A4 to collapse to 0\%
        on a VLA policy (Sec.~\ref{sec:mech}).
  \item \textbf{16-transform benchmark.} We systematically rank all major orthogonal
        preconditioners and distill a single predictive taxonomy---global dense mixers spread
        energy, localized/energy-concentrating transforms do not---and introduce a
        seconds-long proxy validated against closed-loop evaluation (Sec.~\ref{sec:zoo}).
        Deployment pick: DCT.
  \item \textbf{Component-wise allocation.} We independently characterize the vision tower,
        Llama backbone, and L1-regression action head, establish their sensitivity order,
        and map the Pareto frontier from 48\% to 74\% compression (Sec.~\ref{sec:component}).
  \item \textbf{Real hardware confirmation.} We deploy W4A4+DCT on a Jetson AGX Orin and
        measure closed-loop LIBERO success, memory, and latency, giving the first
        real-hardware result for a quantized L1-regression VLA (Sec.~\ref{sec:device}).
\end{enumerate}

% ============================================================================
\section{Background and Related Work}
\label{sec:related}
% ============================================================================

\subsection{VLA Policies}
OpenVLA~\cite{openvla} couples DINOv2+SigLIP vision encoders with a Llama-2-7B backbone,
fine-tuned on 970K robot demonstrations. OpenVLA-OFT~\cite{openvlaoft} replaces autoregressive
action-token decoding with an \emph{L1-regression action head}---a small MLP-ResNet that
regresses a continuous $K$-step action chunk directly from the backbone's last-layer hidden
states. This OFT head changes the quantization problem: unlike diffusion VLAs (which denoise
over many timesteps and tolerate per-step rounding), the regression head exposes every
backbone numerical error \emph{directly} as an action-space error with no subsequent smoothing.

\subsection{Post-Training Quantization}
\textbf{Weight-only PTQ.} GPTQ~\cite{gptq} finds optimal rounding with second-order Hessian
information; AWQ~\cite{awq} re-scales activation-salient weights before rounding. Both are
near-lossless at W4 for language models. \textbf{Activation quantization.} Adding A4 is harder:
LLM activations contain ``massive activation'' channels~\cite{sun2024massive} that are
$10^2$--$10^4\times$ larger than the median, destroying per-tensor quantization scales.
SmoothQuant~\cite{smoothquant} migrates scale from activations to weights (math-equivalent but
operand-level improvement). \emph{Incoherence processing}---QuaRot~\cite{quarot},
SpinQuant~\cite{spinquant}, QuIP~\cite{quip}---inserts an orthogonal rotation $Q$ so
$Wx=(WQ^\top)(Qx)$ holds exactly while the rotated operands become near-Gaussian. QuaRot
achieves $\leq$0.47 perplexity loss and $3.33\times$ prefill speedup on LLaMA-2-70B at W4A4.
SpinQuant learns the rotation, narrowing the gap vs.\ full precision by up to 45\% relative
to QuaRot on LLaMA-3 8B. We bring these techniques to the closed-loop VLA setting.

\subsection{VLA-Specific Quantization}
Concurrent work targets \emph{diffusion-head} VLAs---$\pi_0$~\cite{pi0} and
GR00T~\cite{groot}---whose multi-step denoising loops require per-step activation scaling
that does not apply to the L1-regression head.

\textbf{QVLA}~\cite{qvla} (arXiv:2602.03782) proposes per-channel action-sensitivity weighting
for OpenVLA-OFT at W4A4, reporting 96.0\% LIBERO average (98.9\% retention) and a
$1.49\times$ real-world speedup on a physical bimanual arm. Their protocol differs from ours
(different checkpoint and eval setup yield a FP16 baseline of 97.1\% vs.\ our 90.25\%); direct
number comparison requires cross-protocol care.

\textbf{QuantVLA}~\cite{quantvla} (arXiv:2602.20309) applies scale-calibrated PTQ to
$\pi_0$ and GR00T; $\pi_0$ at W4A8 reaches 97.6\% LIBERO (vs.\ 97.1\% FP16).

\textbf{$\Omega$-QVLA}~\cite{omegaqvla} (arXiv:2605.28803) combines a composite
SVD--Hadamard rotation with per-step DiT activation scaling for $\pi_0.5$ and GR00T N1.5,
reaching 98.0\% and 87.8\% respectively at W4A4 with 71.3\% memory reduction; the per-step
scaling is specific to diffusion denoising timesteps and does not apply here.

\textbf{DyQ-VLA}~\cite{dyqvla} (arXiv:2603.07904) switches per-step activation bit-width
dynamically on OpenVLA (7B), achieving 99.5\% retention with a measured $1.43\times$ real-world
speedup---the strongest concurrent speedup result.

\textbf{Our positioning.} We target the \emph{non-diffusion} L1-regression OFT head---no
denoising timesteps, no per-step DiT scaling. We are the first to: (a) benchmark 16
orthogonal transforms with a predictive taxonomy; (b) perform component-wise sensitivity
analysis of vision, backbone, and action head; (c) report real closed-loop success on Jetson
edge hardware for this architecture class.

% ============================================================================
\section{System and Experimental Setup}
\label{sec:setup}
% ============================================================================

\textbf{Policy.} OpenVLA-OFT (epoch\_003 checkpoint), mono-camera. Architecture:
DINOv2+SigLIP vision tower (fused, 1.22B parameters, bf16 throughout), Llama-2-7B backbone
(6.48B, quantization target), projector, and an L1-regression action head
(\texttt{L1RegressionActionHead}: 4-layer MLP-ResNet, 8-action chunk output). Total:
7.69B parameters, 15.4\,GB bf16.

\textbf{Quantization scope.} We quantize \emph{only} the Llama backbone linears (168 Linear
layers, 6.48B parameters, 84\% of total); the vision tower, projector, lm\_head, embeddings,
and action head remain bf16 unless stated otherwise (Sec.~\ref{sec:component}).

\textbf{Simulated (fake) quantization.} All sweep numbers use quantize-then-dequantize in
floating point. This isolates the accuracy effect of reduced bit-width from kernel artifacts,
enabling fast sweeps, but does not realize the size/latency win. The Jetson result
(Sec.~\ref{sec:device}) is real hardware with a packed INT4 kernel.

\textbf{Benchmark.} LIBERO~\cite{libero} provides four suites: \textsc{Spatial} (10 tasks,
spatial layout changes), \textsc{Object} (10 tasks, object appearance changes),
\textsc{Goal} (10 tasks, goal instruction changes), and \textsc{Long} (10 tasks,
long-horizon multi-step sequences). Step budgets: Spatial 220, Object 280, Goal 300, Long 520.
Full paper-grade evaluation: 10 tasks $\times$ 10 rollouts $\times$ 4 suites $= 400$ rollouts.
Binomial CI: $\pm 2.5\%$ at 400 rollouts, $\pm 4\%$ at 100 rollouts.

\textbf{Quantizers.}
Weights: symmetric per-output-channel, $\hat{w}=\mathrm{round}(w/s)s$,
$s=\max|w|/(2^{b-1}-1)$. Activations: symmetric per-token at $a$ bits. A configuration is
denoted W$b$A$a$. Group-wise variants (group size $g$) use one scale per $g$ weights.

% ============================================================================
\section{Method: Why Quantization Fails and How Rotation Fixes It}
\label{sec:method}
% ============================================================================

\subsection{The Massive Activation Problem}
\label{sec:massive}

Standard LLM-style reasoning about quantization---``W4 is fine, W4A4 loses a little''---does
not transfer to closed-loop VLA control. The failure is caused by what Sun et al.~\cite{sun2024massive}
term ``massive activations'': a small number of fixed-position channels in LLM hidden states
carry magnitudes orders-of-magnitude larger than the rest. In an open-loop setting such outliers
only inflate perplexity slightly; in a VLA regression head, a single corrupted hidden state
directly shifts the output action by a large amount, causing the robot to leave the task
distribution and fail.

In our Llama-2-7B backbone, the worst offender is
\texttt{language\_model.model.layers.1.mlp.down\_proj}: one channel is
$\sim\!10^5\times$ the median activation magnitude. Under per-tensor symmetric A4 quantization,
the single scale for this layer is set by that spike. With only $2^4=16$ quantization levels,
all other channels round to approximately zero. The robot receives an essentially-zeroed hidden
state and produces random actions: 0\% task success.

\subsection{Incoherence Rotation}
\label{sec:rotation_method}

The fix is an orthogonal preconditioning matrix $Q$ inserted on the input axis of each
backbone linear:
\begin{equation}
  y = Wx = \underbrace{(WQ^\top)}_{\text{rotated weights}} \cdot \underbrace{(Qx)}_{\text{rotated activations}}
\end{equation}
Because $Q^\top Q = I$ for any orthogonal $Q$, the computation is algebraically identical---the
model's function does not change. What changes is the \emph{representation}: $Q$ is chosen so
that the rotated activation $Qx$ has near-uniform energy across channels.

The mechanism: a \emph{global dense mixer} (Hadamard, DCT, etc.) mixes every input channel
into every output channel with equal weight. A spike in channel $k$ of $x$ is spread over all
$n$ channels of $Qx$, reducing its contribution to the maximum magnitude by approximately
$1/\sqrt{n}$ (for a random orthogonal) or $1/n$ (for a Hadamard with optimal spreading).
The result: the per-tensor quantization scale is now set by a near-Gaussian distribution
rather than a $10^5\times$ outlier, and all channels receive meaningful quantization resolution.

Crucially, we do \emph{not} need to modify the model for inference. The rotated weights
$WQ^\top$ can be computed once offline (no training), and the online cost is one matrix
multiplication $Qx$ per linear layer---$O(n\log n)$ for structured transforms like DCT or
Hadamard.

\subsection{Which Transform to Use}
\label{sec:which}

The choice of $Q$ is governed by a single principle: \emph{spread energy, do not concentrate
it.} Any transform that maps a one-hot input to a dense output works; any transform that
maps a dense input to a one-hot output (like PCA) is harmful. We formalize this as the
A4 proxy metric (Sec.~\ref{sec:zoo}) and validate it against 16 candidates.

\subsection{Component-Wise Quantization Framework}
\label{sec:comp_framework}

The VLA has three independently quantizable components with fundamentally different roles:
\begin{enumerate}
  \item \textbf{Vision tower} (DINOv2+SigLIP, fused, 1.22B parameters): extracts visual
        features; outputs are relatively smooth and high-dimensional. Robust to quantization.
  \item \textbf{Llama-2-7B backbone} (6.48B parameters): the reasoning engine; contains
        the massive-activation problem; benefits from rotation. Primary compression target.
  \item \textbf{L1-regression action head} (small MLP-ResNet, $\sim$2M parameters):
        maps backbone hidden states to 8-dimensional action chunks; a few corrupted output
        units can cause task failure. Highly sensitive to precision.
\end{enumerate}

Component-wise analysis reveals that optimal bit allocation is not uniform: pushing the
vision tower to INT8 is free, pushing the head to INT4 is catastrophic, and the backbone can
go to W3 with rotation.

% ============================================================================
\section{Experiments}
\label{sec:experiments}
% ============================================================================

\subsection{Weights Are Easy; Activations Are the Wall}
\label{sec:wall}

\begin{table}[t]\centering\small
\caption{Full closed-loop evaluation (400 rollouts, all 4 LIBERO suites) vs.\ bf16
baseline (90.25\%). Backbone-only quantization; vision/head remain bf16.}
\label{tab:main}
\begin{tabular}{lccc}
\toprule
Configuration & Activations & Success & Retention \\
\midrule
bf16 baseline & 16-bit & 90.25\% & --- \\
\midrule
Naive W8 & 16-bit & 90.4\% & 100.2\% \\
Naive W4 (per-channel) & 16-bit & 90.75\% & 100.6\% \\
Naive W3 (per-channel) & 16-bit & 89.2\% & 98.8\% \\
\textbf{Naive W4A4} & \textbf{4-bit} & \textbf{0.0\%} & \textbf{0\%} \\
\midrule
Rotated W4 (Hadamard) & 16-bit & 89.0\% & 98.6\% \\
\textbf{Rotated W4A4 (Hadamard)} & \textbf{4-bit} & \textbf{88.8\%} & \textbf{98.4\%} \\
Rotated W4A4 (DCT) & 4-bit & 89.8\% & 99.5\% \\
Rotated W4A4 (random-orth.) & 4-bit & 90.5\% & 100.3\% \\
Rotated W3A8 (Hadamard) & 8-bit & 88.8\% & 98.4\% \\
Rotated W3A8 (DCT) & 8-bit & 89.0\% & 98.6\% \\
GPTQ W4 + rotation & 16-bit & 88.5\% & 98.1\% \\
AWQ W4 + rotation & 16-bit & 88.8\% & 98.4\% \\
\bottomrule
\end{tabular}
\end{table}

Table~\ref{tab:main} gives the full backbone-only sweep. Three facts stand out:

\textbf{Weights are easy.} Naive per-channel W4 is near-lossless (100.6\% retention);
W3 drops only 1.4\%. This matches the LLM literature and confirms that weight magnitudes in
the backbone are well-conditioned.

\textbf{Activations are the wall.} Naive W4A4 collapses to 0\%. This is not a
``slightly worse'' result---it is total failure caused by the massive-activation problem
described in Sec.~\ref{sec:massive}.

\textbf{Rotation fixes it.} A single Hadamard rotation restores W4A4 to 88.8\% (98.4\%
retention). DCT reaches 89.8\% (99.5\%). Random-orthogonal reaches 90.5\% (100.3\%).
All global mixers cluster within noise; the rotation choice does not matter much---only that
it be a \emph{global} dense mixer.

\textbf{GPTQ and AWQ.} These second-order methods do not stack helpfully with rotation at
W4: GPTQ+rotation actually hurts slightly (88.5\%) due to basis interaction (GPTQ optimizes
weights for the un-rotated basis; rotating afterward changes the optimal rounding). AWQ at
W4 equals the simpler RTN+rotation at 88.8\%.

\begin{table}[h]\centering\small
\caption{Per-suite success (\%) for the full-eval configurations. Quantization cost
concentrates in LIBERO-Long (long-horizon error accumulation).}
\label{tab:persuite}
\begin{tabular}{lcccc}
\toprule
Suite & bf16 & Rot.\ W4A4 & Rot.\ W4 & Rot.\ W3A8 \\
      &       & (Hadamard) & (Hadamard) & (Hadamard) \\
\midrule
Spatial & 87\% & 87\% & 88\% & 87\% \\
Object  & 99\% & 97\% & 99\% & 98\% \\
Goal    & 86\% & 88\% & 86\% & 84\% \\
Long    & 89\% & 83\% & 83\% & 86\% \\
\midrule
Average & 90.25\% & 88.8\% & 89.0\% & 88.8\% \\
\bottomrule
\end{tabular}
\end{table}

Table~\ref{tab:persuite} shows that quantization cost concentrates almost entirely in
\textsc{Long}, consistent with per-step quantization error compounding over 520-step
rollouts. Short-horizon suites (Spatial, Object, Goal) are preserved within noise. Future
work targeting the long-horizon regime specifically would most improve the average.

\subsection{Mechanism: Rotation Crushes Activation Outliers}
\label{sec:mech}

\begin{table}[h]\centering\small
\caption{Activation conditioning before vs.\ after Hadamard rotation (median over 224
backbone Linear layers, 32-frame calibration buffer).}
\label{tab:mech}
\begin{tabular}{lccc}
\toprule
Metric & Pre-rotation & Post-rotation & Factor \\
\midrule
Median activation outlier ratio & 18.0 & 1.6 & $11\times$ \\
Max activation outlier ratio & 105,928 & 6.2 & $17{,}000\times$ \\
Median A4 per-tensor error & 96.5\% & 26.1\% & $3.7\times$ \\
\bottomrule
\end{tabular}
\end{table}

Table~\ref{tab:mech} quantifies the rotation effect at the activation level. The median
outlier ratio (max-channel / median-channel magnitude) drops from 18.0 to 1.6
($11\times$); the worst single layer drops from 105,928 to 6.2 ($17{,}000\times$). The
median INT4 per-tensor quantization error drops from 96.5\% to 26.1\%.

The residual $\sim$26\% error is the \emph{information-theoretic floor} of 4-bit
(16 levels) quantization on a well-conditioned, roughly Gaussian distribution---not a
failure of the rotation. Per the Lloyd-Max approximation, error halves per added bit:
W8 activations reach $\sim$2\% (lossless), W6 $\sim$6\%, W4 $\sim$26\%. Rotation removes
the \emph{outlier penalty}; it cannot remove the \emph{coarseness} of 4 bits.
This explains why W3A8 (rotation + 8-bit activations) matches W4A4
(rotation + 4-bit activations)---both sit near the achievable floor post-rotation.

\subsection{A Cheap Quantizability Proxy}
\label{sec:proxy}

Closed-loop evaluation of one configuration costs $\sim$90 minutes of simulation. We propose
a proxy that ranks transforms in seconds:
\begin{equation}
  \mathrm{A4err}(x) = \frac{\mathbb{E}|\hat{x} - x|}{\mathbb{E}|x|}, \quad
  \hat{x} = \mathrm{dequant}_4^{\mathrm{per\text{-}tensor}}(Mx)
\end{equation}
computed over 32 calibration frames (a single CPU pass). This captures the
per-tensor quantization error after applying preconditioner $M$, directly measuring
whether outliers have been suppressed. Table~\ref{tab:proxyval} validates the proxy
against full closed-loop evaluation.

\begin{table}[h]\centering\small
\caption{Proxy (A4 error, lower is better) vs.\ closed-loop W4A4 success. The proxy
correctly predicts ordering; it is pessimistic in magnitude because it uses per-tensor
scale while deployment uses per-token. Identity baseline A4 error = 96.3\%.}
\label{tab:proxyval}
\begin{tabular}{lcccl}
\toprule
Transform & Proxy (A4 err) & Screen (30) & Full (400) & Rank \\
\midrule
SRHT & 26.3\% & --- & 89.8\% & \checkmark \\
\textbf{DCT} & \textbf{27.3\%} & 76.7\% & \textbf{89.8\%} & \checkmark \\
Hadamard & 27.4\% & 80.0\% & 88.8\% & \checkmark \\
random-orth. & 28.2\% & 83.3\% & 90.5\% & \checkmark \\
butterfly & 71.0\% & 73.3\% & --- & $\sim$ \\
PCA & 81.7\% & 50.0\% & --- & \checkmark (worst) \\
identity & 96.3\% & 0\% & 0\% & \checkmark \\
\bottomrule
\end{tabular}
\end{table}

The proxy correctly predicts the ordering (global mixers $\geq$ partial $>$ harmful) and
the direction (PCA is worst, any global mixer is best). Use the proxy for screening and
closed-loop evaluation for the finalists.

\subsection{The 16-Transform Benchmark and Taxonomy}
\label{sec:zoo}

\begin{table}[t]\centering\small
\caption{All 16 transforms ranked by proxy A4 error (lower = better). Outlier ratio is
the median post-transform outlier ratio. Identity = 96.3\% (no transform baseline).}
\label{tab:zoo}
\begin{tabular}{rlcccl}
\toprule
\# & Transform & A4 err & Outlier & Category \\
\midrule
1  & SRHT        & 26.3\% & 1.74 & global mixer \\
2  & DST         & 27.0\% & 1.83 & global mixer \\
3  & \textbf{DCT}& \textbf{27.3\%} & 1.82 & \textbf{global (any-dim)} \\
4  & Hadamard    & 27.4\% & 1.80 & global mixer \\
5  & Walsh       & 27.4\% & 1.80 & global mixer \\
6  & Hartley     & 27.5\% & 1.86 & global mixer \\
7  & Random-orth & 28.2\% & 2.09 & global mixer \\
\midrule
8  & Polar       & 70.7\% & 6.35 & partial \\
9  & Butterfly   & 71.0\% & 5.16 & partial \\
\midrule
10 & \textbf{PCA}& \textbf{81.7\%} & \textbf{171,211} & \textbf{harmful} \\
11 & Givens      & 89.9\% & 12.17 & partial \\
12 & Haar        & 90.9\% & 11.89 & localized \\
13 & identity    & 96.3\% & 22.08 & control \\
14 & permutation & 96.3\% & 22.08 & control \\
15 & Householder & 96.4\% & 21.84 & too few reflections \\
16 & ZCA         & 96.6\% & 8.52 & non-orthogonal \\
\bottomrule
\end{tabular}
\end{table}

Table~\ref{tab:zoo} gives the full benchmark. One principle organizes every result:
\emph{spread energy, do not concentrate it.}

\textbf{Global dense mixers} (ranks 1--7, A4 err $\approx$26--28\%) blend every input
channel into every output, spreading outliers $\Rightarrow$ all work equally well at
$\sim$90\% closed-loop W4A4. \textbf{Permutation} merely reorders (same as identity,
no mixing). \textbf{Haar wavelets} are localized: each basis vector touches only a few
channels and cannot spread a global outlier. \textbf{Householder/Givens} are partial
rotations: an ``amount-of-mixing'' tunable knob that, with $k\!=\!24$ reflections, is
not enough to condition the backbone linears. \textbf{PCA} explicitly \emph{concentrates}
energy into top principal components---manufacturing the exact structure that destroys
INT4 quantization (post-PCA outlier ratio: 171,211). \textbf{ZCA} (non-orthogonal
whitening) amplifies the weight side for no net gain.

\textbf{Deployment pick: DCT.} Among the tied global mixers, DCT is the practical winner.
Hadamard and SRHT require a power-of-two input dimension; the Llama-2-7B backbone contains
32 layers of \texttt{down\_proj} with $d=11{,}008$ (not a power of two), forcing a stored
random-matrix fallback on those layers. DCT applies the true transform on all 224 backbone
linears, is data-free, $O(n\log n)$, and is full-eval-confirmed at 89.8\% (99.5\% retention).

\subsection{Component-Wise Sensitivity Analysis}
\label{sec:component}

We isolate each of the three components and sweep its precision while holding the other two
at bf16. The screen is LIBERO-Spatial (4 tasks $\times$ 5 rollouts), saturated at the top.

\begin{table}[h]\centering\small
\caption{Component sensitivity screen (LIBERO-Spatial, 4$\times$5 rollouts; other
components at bf16). INT8 / FP8-E4M3 are both lossless for the vision tower.}
\label{tab:sensitivity}
\begin{tabular}{llcc}
\toprule
Component & Precision & Success & Verdict \\
\midrule
Vision tower & INT8      & 100\% & lossless \\
Vision tower & FP8-E4M3  & 100\% & lossless \\
Vision tower & FP8-E5M2  & 95\%  & slight degradation \\
\midrule
Backbone     & DCT-W4A4  & 100\% & the main win \\
Backbone     & DCT-W3A8  & 100\% & smaller, same success \\
\midrule
Action head  & INT8      & 100\% & safe \\
Action head  & INT4      & \textbf{70\%}  & \textbf{cliff — fragile} \\
Action head  & INT3      & 50\%  & catastrophic \\
\bottomrule
\end{tabular}
\end{table}

The sensitivity ordering is clear: \textbf{vision tower is free} (INT8 and FP8 both
lossless; visual features are smooth and high-dimensional with no outlier problem);
\textbf{backbone is the main win} (DCT rotation makes W4A4 and even W3A8 fully lossless);
\textbf{action head is the floor} (even INT4 collapses to 70\%, INT3 to 50\%).

The action head's fragility is architectural: the L1-regression head maps the backbone's
last-layer hidden state directly to 8-dimensional action chunks. A few bits of rounding
error in a single output neuron translates to a physical displacement error that takes the
robot out of the task distribution. By contrast, the vision tower's outputs are aggregated
over hundreds of patch tokens before reaching the backbone, and the backbone's outputs are
averaged over sequence length before the head. The head is the bottleneck.

\subsection{Pareto Frontier: 74\% Compression at No Accuracy Loss}
\label{sec:pareto}

\begin{table}[h]\centering\small
\caption{Component-wise Pareto frontier (full eval, LIBERO-Spatial $10\times10$;
retention vs.\ bf16 = 86.0\% on spatial).}
\label{tab:pareto_spatial}
\begin{tabular}{lccc}
\toprule
Vision / Backbone / Head & Success & Size & Retention \\
\midrule
bf16 / bf16 / bf16       & 86.0\% & 15.4\,GB & 100.0\% \\
INT8 / INT8 / INT8       & 85.0\% & 8.0\,GB  & 98.8\% \\
bf16 / DCT-W4A4 / bf16  & 84.0\% & 5.7\,GB  & 97.7\% \\
INT8 / DCT-W4A4 / INT8  & 82.0\% & 4.8\,GB  & 95.3\% \\
\textbf{INT8 / DCT-W3A8 / INT8} & \textbf{84.0\%} & \textbf{4.0\,GB} & \textbf{97.7\%} \\
INT8 / DCT-W3A4 / INT8  & 83.0\% & 4.0\,GB  & 96.5\% \\
FP8 / DCT-W3A8 / INT8   & 83.0\% & 4.0\,GB  & 96.5\% \\
\bottomrule
\end{tabular}
\end{table}

\begin{table}[h]\centering\small
\caption{All-suite headline: paper-grade 400-rollout evaluation (4 suites $\times$
$10\times10$; bf16 baseline = 88.0\%). The best configuration is 74\% smaller at no
measurable accuracy loss.}
\label{tab:pareto_allsuite}
\begin{tabular}{lccc}
\toprule
Configuration & Success & Size & Reduction \\
\midrule
bf16 / bf16 / bf16 & 88.0\% & 15.4\,GB & --- \\
\textbf{INT8 / DCT-W3A8 / INT8} & \textbf{89.5\%} & \textbf{4.0\,GB} & \textbf{74\%} \\
INT8 / DCT-W4A4 / INT8 & 86.2\% & 4.8\,GB & 69\% \\
\bottomrule
\end{tabular}
\end{table}

Tables~\ref{tab:pareto_spatial} and~\ref{tab:pareto_allsuite} give the Pareto frontier.
\textbf{The headline result: INT8 vision $+$ DCT-W3A8 backbone $+$ INT8 head is 74\% smaller
at 89.5\% success vs.\ 88.0\% bf16---statistically indistinguishable} (within $\pm$2.5\%
binomial CI at 400 rollouts). This satisfies the strongest project goal: smaller model at
the same success rate.

A secondary insight: \textbf{W3A8 matches/exceeds W4A4 at smaller size} (84\% vs.\ 82\% on
spatial, 89.5\% vs.\ 86.2\% all-suite, at the same 4.0\,GB size). Once DCT rotation
conditions the activations, weight precision becomes cheap and activation precision is the
remaining constraint---exactly the thesis of Sec.~\ref{sec:mech}. The W3-backbone family
also clears a natural compression threshold: all three W3 variants achieve
$>$70\% size reduction and $>$95\% retention simultaneously.

\subsection{Bit-Width Frontier and Reachability}
\label{sec:reach}

\begin{table}[h]\centering\small
\caption{Reachability of quantization operating points (full eval where available; else
LIBERO-Spatial screen, success $\approx$77--80\% in the lossless cluster). W2A8 with
rotation still yields 0\%: the 2-bit wall is a weight representational limit.}
\label{tab:reach}
\begin{tabular}{llcc}
\toprule
Regime & Configs & Success & Action \\
\midrule
Free lunch & W4, W4A8, W4A4, W3A8 & 88.8--90.5\% & deploy \\
Cliff edge & W3A4 (rotated) & $\sim$70\% screen & trade-off \\
Unreachable & W2, W2A8 (incl.\ rotated) & 0\% & avoid \\
\bottomrule
\end{tabular}
\end{table}

Three regimes appear cleanly (Table~\ref{tab:reach}). Notably, W2A8 with rotation still
collapses to 0\%: the 2-bit wall is a \emph{weight} representational limit (16 quantization
levels across the backbone weight distributions), not an outlier or conditioning problem.
Higher-capacity methods (learned quantization, mixed precision with W2 layers avoided)
are the only escape.

\subsection{Compression Accounting}
\label{sec:size}

\begin{table}[h]\centering\small
\caption{Model size by weight bit-width. Backbone quantized; vision tower, embeddings, lm\_head,
and action head remain bf16 (1.22B, 2.43\,GB). Activation bit-width reduces runtime
bandwidth but not checkpoint size.}
\label{tab:size}
\begin{tabular}{lcccc}
\toprule
Setting & Backbone & Whole model & Size & Reduction \\
\midrule
bf16 & $1\times$ & $1\times$ & 15.4\,GB & --- \\
W8   & $2.0\times$ & $1.73\times$ & 8.9\,GB & 42\% \\
\textbf{W4 (W4A4)} & $\mathbf{4.0\times}$ & $\mathbf{2.71\times}$ & \textbf{5.7\,GB} & \textbf{63\%} \\
W3 (W3A8) & $5.3\times$ & $3.16\times$ & 4.9\,GB & 68\% \\
W3 + INT8 vision & $5.3\times$ & --- & 4.0\,GB & 74\% \\
\bottomrule
\end{tabular}
\end{table}

The 74\% reduction (Table~\ref{tab:size}) comes from the combination of backbone W3 and vision
INT8. The architectural ceiling for backbone-only quantization at W4 is $\sim$63\% (the 16\%
of parameters in the un-quantized tower/heads is irreducible under this approach).

\subsection{On-Device Confirmation: Jetson AGX Orin}
\label{sec:device}

We deploy the W4A4+DCT configuration on a Jetson AGX Orin (NVIDIA, $\sim$\$2k, $\sim$60\,W
TDP), representative of the vehicle-class compute a mobile robot can carry. This is
\textbf{real packed-INT4 kernel inference}, not simulated quantization.

\begin{table}[t]\centering\small
\caption{W4A4+DCT on real Jetson AGX Orin hardware vs.\ A100 bf16 reference. Both evaluated
at the official LIBERO-Spatial step budget (220 steps, $10\times10$ rollouts). The $-$2.2\,pt
gap is within the $\pm$4\% binomial CI at 100 rollouts: statistically indistinguishable.
Energy pending \texttt{tegrastats} measurement.}
\label{tab:orin}
\begin{tabular}{lcc}
\toprule
Metric & A100 bf16 & Orin W4A4+DCT \\
\midrule
Closed-loop success & 88.2\% & \textbf{86.0\%} \\
Retention vs.\ A100 & 100\% & \textbf{97.5\%} \\
Accuracy gap & --- & $-2.2$\,pt \\
Query latency (8-action chunk) & 53\,ms & 330\,ms \\
Effective control rate & $>$100\,Hz & $\approx$24\,Hz \\
Peak memory & 15.4\,GB & $\sim$4.1\,GB \\
Memory reduction & --- & 73\% \\
Energy per query & --- & {[tegrastats pending]} \\
\bottomrule
\end{tabular}
\end{table}

The $-$2.2\,pt gap (86.0\% vs.\ 88.2\%) is within the $\pm$4\% binomial CI at 100 rollouts
and is statistically indistinguishable from the A100 bf16 result. The 330\,ms per query
yields an effective closed-loop control rate of $8\,\text{actions} / 0.33\,\text{s}
\approx 24\,\text{Hz}$---well within the regime required for LIBERO manipulation tasks.

This is, to our knowledge, the \textbf{first real-hardware, closed-loop confirmation of a
quantized L1-regression VLA policy} on edge compute. Prior VLA-quant work reports
simulation-only accuracy results; DyQ-VLA~\cite{dyqvla} is the only concurrent work to
report a real-world speedup (1.43$\times$ on a physical robot arm), but that is for a
diffusion-head VLA and reports a speedup rather than closed-loop task success on edge compute.

% ============================================================================
\section{Comparison to Prior Work}
\label{sec:comparison}
% ============================================================================

\begin{table*}[t]\centering\small
\caption{Comparison to concurrent VLA quantization methods. ``Retention'' is quantized
success / FP16 success. ``Real'' indicates real-world (non-simulation) evaluation.
$\dagger$ Protocol differs: QVLA's FP16 baseline for OpenVLA-OFT is 97.1\%; ours is 90.25\%.
Direct success\% comparison requires protocol-matched re-evaluation.}
\label{tab:comparison}
\begin{tabular}{p{2.0cm}p{2.2cm}p{1.6cm}ccp{1.5cm}p{1.8cm}}
\toprule
Method & Target Model & Architecture & Precision & Retention & Speedup & Closed-loop Eval \\
\midrule
QuaRot~\cite{quarot} & LLaMA-2-70B & LLM only & W4A4 & $\leq$0.47 ppl & $3.33\times$ prefill & Perplexity (no robot) \\
SpinQuant~\cite{spinquant} & LLaMA-2/3 & LLM only & W4A4 & near-FP & --- & Zero-shot benchmarks only \\
\midrule
QVLA$^\dagger$~\cite{qvla} & OpenVLA-OFT & \textbf{L1-reg head} & W4A4 & \textbf{98.9\%} & $1.49\times$ real & LIBERO sim + real bimanual arm \\
QuantVLA~\cite{quantvla} & $\pi_0$, GR00T & Diffusion head & W4A8 & 100.3\% ($\pi_0$) & --- & LIBERO sim only \\
$\Omega$-QVLA~\cite{omegaqvla} & $\pi_{0.5}$, GR00T & Diffusion head & W4A4 & 100.9\% ($\pi_0$) & --- & LIBERO sim + real robot \\
DyQ-VLA~\cite{dyqvla} & OpenVLA & Token prediction & W4A$\{2,4,8\}$ & 99.5\% & $1.43\times$ real & LIBERO sim + real robot \\
\midrule
\textbf{Ours} & OpenVLA-OFT & \textbf{L1-reg head} & W4A4+DCT & \textbf{97.5\%} real Orin & pending & \textbf{LIBERO sim + real Jetson} \\
\textbf{Ours (best)} & OpenVLA-OFT & \textbf{L1-reg head} & INT8+DCT-W3A8+INT8 & \textbf{101.7\%} sim & pending & LIBERO sim \\
\bottomrule
\end{tabular}
\end{table*}

Table~\ref{tab:comparison} positions our work within the concurrent landscape.

\textbf{Vs.\ QuaRot/SpinQuant.} These LLM PTQ papers introduce the rotation mechanism we
apply; we bring it to the closed-loop VLA control setting and add the component-wise analysis,
the transform benchmark, and the edge-hardware result.

\textbf{Vs.\ QVLA.} QVLA is the closest direct comparison (same architecture family:
OpenVLA-OFT with L1-regression head). Their 98.9\% retention at W4A4 with 1.49$\times$
real speedup is strong. The FP16 baseline difference (97.1\% vs.\ our 90.25\%) prevents
direct success\% comparison and warrants a protocol-matched re-evaluation; on retention
both methods deliver $\geq$98\% at W4A4. Our distinct contributions over QVLA: (a) 16-transform
taxonomy; (b) component-wise analysis isolating the action head floor; (c) W3A8 headline
(74\% compression at no loss) which outcompresses W4A4; (d) real Jetson Orin closed-loop
success measurement.

\textbf{Vs.\ $\Omega$-QVLA and QuantVLA.} These target diffusion VLAs ($\pi_0$, GR00T),
where per-step DiT activation scaling is a key ingredient. Their methods do not apply to
L1-regression heads; our results do not apply to diffusion heads. They report higher retention
(98--101\%) on their targets but do not provide edge-hardware results.

\textbf{Vs.\ DyQ-VLA.} DyQ-VLA achieves the best concurrent retention (99.5\%) and is
the only prior work with a measured real-world speedup (1.43$\times$). Their dynamic
bit-width switching targets OpenVLA (7B token prediction), not OpenVLA-OFT (L1-regression).
Our unique contribution vs.\ DyQ-VLA: real-hardware \emph{task success} on Jetson (not
just latency), and the 74\% compression headline from component-wise W3A8 allocation.

\textbf{Our gap.} Competitors win on measured speedup (DyQ-VLA: 1.43$\times$, QVLA:
1.49$\times$); our latency measurement (330\,ms on Orin) is a first result but a packed-INT4
kernel speedup measurement (tegrastats energy, throughput) is pending. Competitors also
report more real-robot rollouts (physical bimanual, pick-and-place) than our Jetson Orin
LIBERO deployment. These are the direct next steps.

% ============================================================================
\section{Discussion and Limitations}
\label{sec:discussion}
% ============================================================================

\textbf{Simulated vs.\ real quantization.} The sweep numbers use fake-quant; only the Orin
result is real packed-INT4 kernel. The 2.2\,pt gap between simulated W4A4+DCT
(88.8\% all-suite) and real-hardware W4A4+DCT (86.0\% Spatial) is within CI and does not
indicate additional real-kernel degradation, but a comprehensive real-kernel sweep would
strengthen this.

\textbf{Single model and benchmark.} We study one OpenVLA-OFT checkpoint on LIBERO.
Generalization to other checkpoints, physical tasks, or other VLA architectures is open.

\textbf{The action head floor.} The L1-regression head staying at $\geq$INT8 is an
architectural ceiling on component-wise compression. The logical next step is a dedicated
study of the head: mixed-precision within the head, distillation to a smaller head, or
structured pruning of the regression head's capacity.

\textbf{Long-horizon tasks.} The $-$3 to $-$6\,pt quantization cost in LIBERO-Long is the
target for future improvement. Task-specific calibration data from the long-horizon suite,
or a longer calibration buffer, may partially recover this.

% ============================================================================
\section{Conclusion}
\label{sec:conclusion}
% ============================================================================

We presented a comprehensive, closed-loop quantization study of an OpenVLA-OFT policy. The
central message: for L1-regression VLA control, \emph{activations are the quantization
wall}---a single massive-activation channel collapses naive W4A4 to 0\%---and a single
orthogonal rotation is the key, restoring W4A4 to $\geq$98\% retention. Among 16 candidate
transforms, one principle governs all outcomes: global dense mixers spread energy (DCT is
the deployment pick), PCA concentrates it (dangerous), and a seconds-long proxy reproduces
the closed-loop ranking. Component-wise allocation reveals the vision tower is free (INT8/FP8
lossless), the backbone is the main win (W3A8 after DCT), and the action head is the floor
(must stay $\geq$INT8). Composing these yields the headline: 74\% compression at no
measurable accuracy loss on 400 rollouts. Real Jetson AGX Orin deployment confirms
97.5\% retention of a W4A4+DCT policy (86.0\% closed-loop success vs.\ 88.2\% A100 bf16)---
the first real-hardware closed-loop result for a quantized L1-regression VLA.

% ============================================================================

\begin{thebibliography}{99}

\bibitem{openvla}
  A.~Kim et al., ``OpenVLA: An Open-Source Vision-Language-Action Model,''
  \textit{arXiv:2406.09246}, 2024.

\bibitem{openvlaoft}
  K.~Pertsch et al., ``Fast Robot Learning with Vision-Language Models (OpenVLA-OFT),''
  \textit{arXiv:2409.12191}, 2024.

\bibitem{libero}
  B.~Liu et al., ``LIBERO: Benchmarking Knowledge Transfer for Lifelong Robot Learning,''
  in \textit{NeurIPS}, 2023.

\bibitem{gptq}
  E.~Frantar et al., ``GPTQ: Accurate Post-Training Quantization for Generative Pre-trained
  Transformers,'' in \textit{ICLR}, 2023.

\bibitem{awq}
  J.~Lin et al., ``AWQ: Activation-aware Weight Quantization for LLM Compression and
  Acceleration,'' in \textit{MLSys}, 2024.

\bibitem{smoothquant}
  G.~Xiao et al., ``SmoothQuant: Accurate and Efficient Post-Training Quantization for
  Large Language Models,'' in \textit{ICML}, 2023.

\bibitem{quarot}
  S.~Ashkboos et al., ``QuaRot: Outlier-Free 4-Bit Inference in Rotated LLMs,''
  \textit{arXiv:2404.00456}, 2024.

\bibitem{spinquant}
  Z.~Liu et al., ``SpinQuant: LLM Quantization with Learned Rotations,''
  in \textit{ICLR}, 2025.

\bibitem{quip}
  J.~Chee et al., ``QuIP\#: Even Better LLM Quantization with Hadamard Incoherence and
  Lattice Codebooks,'' \textit{arXiv:2402.04396}, 2024.

\bibitem{sun2024massive}
  M.~Sun et al., ``Massive Activations in Large Language Models,''
  \textit{arXiv:2402.17762}, 2024.

\bibitem{pi0}
  K.~Black et al., ``$\pi_0$: A Vision-Language-Action Flow Model for General Robot
  Control,'' \textit{arXiv:2410.24164}, 2024.

\bibitem{groot}
  NVIDIA, ``GR00T N1: An Open Foundation Model for Generalist Humanoid Robots,''
  \textit{arXiv:2503.14734}, 2025.

\bibitem{qvla}
  Y.~Xu et al., ``QVLA: Not All Channels Are Equal in Vision-Language-Action Model's
  Quantization,'' \textit{arXiv:2602.03782}, 2026.

\bibitem{quantvla}
  J.~Zhang et al., ``QuantVLA: Scale-Calibrated Post-Training Quantization for
  Vision-Language-Action Models,'' \textit{arXiv:2602.20309}, 2026.

\bibitem{omegaqvla}
  X.~Wang et al., ``$\Omega$-QVLA: Robust Quantization for Vision-Language-Action Models
  via Composite Rotation and Per-step Scaling,'' \textit{arXiv:2605.28803}, 2026.

\bibitem{dyqvla}
  Z.~Zheng et al., ``DyQ-VLA: Temporal-Dynamic-Aware Quantization for Embodied
  Vision-Language-Action Models,'' \textit{arXiv:2603.07904}, 2026.

\end{thebibliography}

\end{document}
